\clearpage
\begin{explanationbox}
\par\medskip
\noindent\begin{minipage}{\linewidth}
\paragraph{从同一高度的 A 点出发，回到 B 点，会经过什么？}
\begin{center}% F081: precise mathematical teaching diagram.
\begin{tikzpicture}[x=1cm,y=1cm,>=stealth,font=\small]
\draw[rounded corners=3pt,gray!40] (0,0) rectangle (14.2,8.5);
\begin{scope}[yshift=1cm]
\node at (7.1,6.9) {两端等高：上拱、平走、下凹};
\node[] at (2.6,6.3) {先上升，再下降};
\begin{scope}[shift={(2.6,3.9)},x=1.35cm,y=1.4cm]
\draw[->,gray!65] (-1.35,0)--(1.35,0) node[right] {$x$};
\draw[->,gray!65] (0,-1.45)--(0,1.2) node[left] {$y$};
\draw[blue!75!black,line width=1.15pt,,domain=-1:1,samples=160] plot (\x,{1-(\x)*(\x)});
\fill[blue!75!black] (-1,0) circle (1.8pt);\node[above left] at (-1,0) {$A$};
\fill[blue!75!black] (1,0) circle (1.8pt);\node[above right] at (1,0) {$B$};
\draw[orange!85!black,line width=1.2pt] (-0.3,1)--(0.3,1);
\fill[orange!85!black] (0,1) circle (1.8pt);
\end{scope}
\node[] at (2.6,1.25) {内部最高处：$f^{\prime}(c)=0$};
\node[] at (7.1,6.3) {全程保持等高};
\begin{scope}[shift={(7.1,3.9)},x=1.35cm,y=1.4cm]
\draw[->,gray!65] (-1.35,0)--(1.35,0) node[right] {$x$};
\draw[->,gray!65] (0,-1.45)--(0,1.2) node[left] {$y$};
\draw[blue!75!black,line width=1.15pt,,domain=-1:1,samples=160] plot (\x,{0});
\fill[blue!75!black] (-1,0) circle (1.8pt);\node[above left] at (-1,0) {$A$};
\fill[blue!75!black] (1,0) circle (1.8pt);\node[above right] at (1,0) {$B$};
\draw[orange!85!black,line width=1.2pt] (-0.3,0)--(0.3,0);
\fill[orange!85!black] (0,0) circle (1.8pt);
\end{scope}
\node[] at (7.1,1.25) {内部每一点：$f^{\prime}(c)=0$};
\node[] at (11.6,6.3) {先下降，再上升};
\begin{scope}[shift={(11.6,3.9)},x=1.35cm,y=1.4cm]
\draw[->,gray!65] (-1.35,0)--(1.35,0) node[right] {$x$};
\draw[->,gray!65] (0,-1.45)--(0,1.2) node[left] {$y$};
\draw[blue!75!black,line width=1.15pt,,domain=-1:1,samples=160] plot (\x,{(\x)*(\x)-1});
\fill[blue!75!black] (-1,0) circle (1.8pt);\node[above left] at (-1,0) {$A$};
\fill[blue!75!black] (1,0) circle (1.8pt);\node[above right] at (1,0) {$B$};
\draw[orange!85!black,line width=1.2pt] (-0.3,-1)--(0.3,-1);
\fill[orange!85!black] (0,-1) circle (1.8pt);
\end{scope}
\node[] at (11.6,1.25) {内部最低处：$f^{\prime}(c)=0$};
\end{scope}
\node[text width=13cm,align=center] at (7.1,.9) {三种典型情形均有水平切线；前提仍是闭区间连续、开区间可导。};
\end{tikzpicture}
\end{center}
\end{minipage}\par\medskip

记 $A(a,f(a))$、$B(b,f(b))$，且 $a<b$、$f(a)=f(b)$。图中有三种最容易想象的走法：
\begin{itemize}
\item \textbf{先上升再下降：}内部有最高处，那里切线水平；
\item \textbf{一直平着走到 B：}内部每一点的导数都是 $0$；
\item \textbf{先下降再上升：}内部有最低处，那里切线水平。
\end{itemize}
这三种图形是\textbf{典型情形，不是说所有函数只能这样单调分成两段}。下面还可以反复起伏。

\textbf{从图像中抓住不变的东西：}起点 $A$ 和终点 $B$ 同高，所以连接它们的 $AB$ 是一条水平线。无论曲线中间怎样起伏，只要符合罗尔定理的条件，\textbf{总能找到一处切线与 $AB$ 平行}。这比只背“导数等于零”更容易想象。

\end{explanationbox}
\clearpage
\begin{explanationbox}
\par\medskip
\noindent\begin{minipage}{\linewidth}
\paragraph{组合起来也成立：起伏再多，至少有一处走平}
\begin{center}% F081: precise mathematical teaching diagram.
\begin{tikzpicture}[x=1cm,y=1cm,>=stealth,font=\small]
\draw[rounded corners=3pt,gray!40] (0,0) rectangle (14.2,8.8);
\begin{scope}[yshift=1cm]
\node at (7.1,7.2) {也可以多次起伏：至少一个，不限制只有一个};
\begin{scope}[shift={(1.4,3.7)},x=2.45cm,y=1.5cm]
\draw[->,gray!65] (-0.2,0)--(4.45,0) node[right] {$x$};
\draw[->,gray!65] (0,-1.45)--(0,1.6) node[above] {$y$};
\draw[blue!75!black,line width=1.15pt,,domain=0:4,samples=160] plot (\x,{sin(180*\x)});
\fill[blue!75!black] (0,0) circle (1.8pt);\node[below left] at (0,0) {$A$};
\fill[blue!75!black] (4,0) circle (1.8pt);\node[below right] at (4,0) {$B$};
\draw[orange!85!black,line width=1.2pt] (0.28,1)--(0.72,1);
\fill[orange!85!black] (0.5,1) circle (1.8pt);
\draw[orange!85!black,line width=1.2pt] (1.28,-1)--(1.72,-1);
\fill[orange!85!black] (1.5,-1) circle (1.8pt);
\draw[orange!85!black,line width=1.2pt] (2.28,1)--(2.72,1);
\fill[orange!85!black] (2.5,1) circle (1.8pt);
\draw[orange!85!black,line width=1.2pt] (3.28,-1)--(3.72,-1);
\fill[orange!85!black] (3.5,-1) circle (1.8pt);
\end{scope}
\node[] at (7.1,1.2) {图例 $f(x)=\sin(\pi x),\quad 0\le x\le4$：有四处水平切线。};
\end{scope}
\node[text width=13cm,align=center] at (7.1,.9) {三种形状可以反复组合；一般曲线不必只有一段上升和一段下降。};
\end{tikzpicture}
\end{center}
\end{minipage}\par\medskip


\paragraph{再看一幅关键的图：等高的两点与中间的水平切线}
\begin{center}% F081 新增图：同高两点与平行的水平切线。保留其他旧图不变。
\begin{tikzpicture}[x=1.7cm,y=1.05cm,>=stealth,font=\small]
  \path[use as bounding box] (-0.45,-0.5) rectangle (6.9,4.7);
  \draw[->,gray!70] (-0.3,0) -- (4.05,0) node[right] {$x$};
  \draw[->,gray!70] (0,-0.25) -- (0,3.95) node[above] {$y$};
  \draw[blue!75!black,very thick,domain=0.4:3.6,samples=120] plot (\x,{(\x-2)^2+1});
  \draw[blue!65,thick,dashed] (0.72,2) -- (4.10,2);
  \draw[orange!85!black,very thick] (1.38,1) -- (4.10,1);
  \fill[blue!75!black] (1,2) circle (2pt) (3,2) circle (2pt);
  \fill[orange!85!black] (2,1) circle (2.3pt);
  \node[above right,xshift=2pt,yshift=8pt] at (1,2) {$A(1,2)$};
  \node[above left,xshift=-2pt,yshift=8pt] at (3,2) {$B(3,2)$};
  \node[below,yshift=-5pt] at (2,1) {$C(2,1)$};
  \node[anchor=west,blue!75!black] at (4.35,2) {$AB: y=2$};
  \node[anchor=west,orange!85!black] at (4.35,1) {水平切线 $y=1$};
  \node[gray!80] at (2,4.30) {$f(x)=(x-2)^2+1$};
\end{tikzpicture}
\end{center}

图中两点 $A(1,2)$、$B(3,2)$ 同高，连接它们得到一条水平线；中间的 $C(2,1)$ 比两端低，但是\textbf{$C$ 处的切线与 $AB$ 一样水平}。因此罗尔定理说的是\textbf{切线方向一样}，不是\textbf{这三个点高度一样}。

\textbf{真正不变的理由：}只要中间不是处处与端点等高，就有一点高于或低于端点。闭区间连续保证函数能取到最高值和最低值；其中至少一个必须在内部。再用这里的可导性，就能证明该处导数为 $0$。

不能只凭图形上的“转弯”就认定导数为零；如果是尖角，导数可能不存在。\textbf{在 $[a,b]$ 上连续、在 $(a,b)$ 内处处可导}，两项缺一不可。

\textbf{两端不一定是零。}例如 $f(a)=f(b)=5$，结论同样成立。把“等高”取成“都为零”，就得到两个不同零点之间必有导数零点。

\textbf{把 $x$ 轴向上或向下平移，也一样能用。}例如一条水平线 $y=h$ 与曲线有三个交点，就把这三个交点按从左到右排好。前两个交点之间至少有一处水平切线，后两个之间也至少有一处；\textbf{因此至少有两处不同的水平切线位置}。$x$ 轴其实只是 $h=0$ 的一种特殊情况。

\textbf{走平不一定就是山顶或谷底。}例如 $f(x)=x^3$ 在 $x=0$ 处切线水平，但图像依然从左向右不断升高。罗尔定理保证的是\textbf{存在水平切线}，不能把它误记成“一定要拐头”。


\textbf{怎样用来数根？}每一对相邻零点都要夹出一个不同的导数零点。因此先数导数的零点，可以限制原方程的解的个数；再用连续性与异号函数值证明解确实存在。这个顺序比硬解高次方程更直接。
\end{explanationbox}


% -------- 新增独立小节：导数为零的九种走向 --------
\clearpage
\begin{explanationbox}
\paragraph{导数为零的九种走向：水平切线不等于曲线掉头}

\textbf{先分清两件事：}在 $x=c$ 处 $f'(c)=0$，只说明\textbf{这一点的切线水平}；
图像到这里是继续上升、掉头下降，还是变成一段平路，\textbf{不能只看这一点}，还要看它\textbf{左边与右边}。

为方便比较，先假设在 $c$ 左右各取一小段，\textbf{每一侧的导数分别始终大于 $0$、始终等于 $0$，或始终小于 $0$}。
于是左侧有三种，右侧也有三种，一共 $3\times3=9$ 种基本组合。
注意：表中的 ``$=0$'' 指\textbf{这一侧的一整小段都等于 $0$}，不是只在 $c$ 这一点等于 $0$。

\begin{center}
\small
\renewcommand{\arraystretch}{1.4}
\begin{tabular}{c|c|c|c}
\textbf{左侧\textbackslash 右侧} & \textbf{$f'>0$（升）} & \textbf{$f'=0$（平）} & \textbf{$f'<0$（降）} \\ \hline
\textbf{$f'>0$（升）} & \textcircled{\raisebox{-0.1ex}{\scalebox{0.72}{1}}} 升 $\to$ 升 & \textcircled{\raisebox{-0.1ex}{\scalebox{0.72}{2}}} 升 $\to$ 平 & \textcircled{\raisebox{-0.1ex}{\scalebox{0.72}{3}}} 升 $\to$ 降 \\ \hline
\textbf{$f'=0$（平）} & \textcircled{\raisebox{-0.1ex}{\scalebox{0.72}{4}}} 平 $\to$ 升 & \textcircled{\raisebox{-0.1ex}{\scalebox{0.72}{5}}} 平 $\to$ 平 & \textcircled{\raisebox{-0.1ex}{\scalebox{0.72}{6}}} 平 $\to$ 降 \\ \hline
\textbf{$f'<0$（降）} & \textcircled{\raisebox{-0.1ex}{\scalebox{0.72}{7}}} 降 $\to$ 升 & \textcircled{\raisebox{-0.1ex}{\scalebox{0.72}{8}}} 降 $\to$ 平 & \textcircled{\raisebox{-0.1ex}{\scalebox{0.72}{9}}} 降 $\to$ 降
\end{tabular}
\end{center}

\begin{center}% F081: 九种导数走向。每个小图在 x=0 处导数都为0。
% 左右分别取 f(x)=a x^3、f(x)=b x^3 (a,b∈{1,0,-1})。
\begin{tikzpicture}[font=\footnotesize, line cap=round, line join=round]
\node[font=\bfseries] at (5.0,9.28) {右侧的走向：上升／平路／下降};
\begin{scope}[shift={(1.75,7.55)},x=.80cm,y=.69cm]
  \draw[gray!35,rounded corners=2pt] (-1.63,-1.98) rectangle (1.63,1.43);
  \draw[gray!55,thin] (-1.28,0)--(1.30,0);
  \draw[gray!55,thin] (0,-1.24)--(0,1.14);
  \draw[blue!70!black,line width=1.2pt,smooth] plot[domain=-1.09:0,samples=26] (\x,{1*\x*\x*\x});
  \draw[blue!70!black,line width=1.2pt,smooth] plot[domain=0:1.09,samples=26] (\x,{1*\x*\x*\x});
  \draw[orange!85!black,thick] (-.34,0)--(.34,0);
  \fill[orange!85!black] (0,0) circle (1.6pt);
  \node[font=\scriptsize] at (0,-1.35) {\textcircled{\raisebox{-0.1ex}{\scalebox{0.72}{1}}} 升→升};
\end{scope}
\begin{scope}[shift={(4.98,7.55)},x=.80cm,y=.69cm]
  \draw[gray!35,rounded corners=2pt] (-1.63,-1.98) rectangle (1.63,1.43);
  \draw[gray!55,thin] (-1.28,0)--(1.30,0);
  \draw[gray!55,thin] (0,-1.24)--(0,1.14);
  \draw[blue!70!black,line width=1.2pt,smooth] plot[domain=-1.09:0,samples=26] (\x,{1*\x*\x*\x});
  \draw[blue!70!black,line width=1.2pt,smooth] plot[domain=0:1.09,samples=26] (\x,{0*\x*\x*\x});
  \draw[orange!85!black,thick] (-.34,0)--(.34,0);
  \fill[orange!85!black] (0,0) circle (1.6pt);
  \node[font=\scriptsize] at (0,-1.35) {\textcircled{\raisebox{-0.1ex}{\scalebox{0.72}{2}}} 升→平};
\end{scope}
\begin{scope}[shift={(8.21,7.55)},x=.80cm,y=.69cm]
  \draw[gray!35,rounded corners=2pt] (-1.63,-1.98) rectangle (1.63,1.43);
  \draw[gray!55,thin] (-1.28,0)--(1.30,0);
  \draw[gray!55,thin] (0,-1.24)--(0,1.14);
  \draw[blue!70!black,line width=1.2pt,smooth] plot[domain=-1.09:0,samples=26] (\x,{1*\x*\x*\x});
  \draw[blue!70!black,line width=1.2pt,smooth] plot[domain=0:1.09,samples=26] (\x,{-1*\x*\x*\x});
  \draw[orange!85!black,thick] (-.34,0)--(.34,0);
  \fill[orange!85!black] (0,0) circle (1.6pt);
  \node[font=\scriptsize] at (0,-1.35) {\textcircled{\raisebox{-0.1ex}{\scalebox{0.72}{3}}} 升→降};
\end{scope}
\begin{scope}[shift={(1.75,5.03)},x=.80cm,y=.69cm]
  \draw[gray!35,rounded corners=2pt] (-1.63,-1.98) rectangle (1.63,1.43);
  \draw[gray!55,thin] (-1.28,0)--(1.30,0);
  \draw[gray!55,thin] (0,-1.24)--(0,1.14);
  \draw[blue!70!black,line width=1.2pt,smooth] plot[domain=-1.09:0,samples=26] (\x,{0*\x*\x*\x});
  \draw[blue!70!black,line width=1.2pt,smooth] plot[domain=0:1.09,samples=26] (\x,{1*\x*\x*\x});
  \draw[orange!85!black,thick] (-.34,0)--(.34,0);
  \fill[orange!85!black] (0,0) circle (1.6pt);
  \node[font=\scriptsize] at (0,-1.35) {\textcircled{\raisebox{-0.1ex}{\scalebox{0.72}{4}}} 平→升};
\end{scope}
\begin{scope}[shift={(4.98,5.03)},x=.80cm,y=.69cm]
  \draw[gray!35,rounded corners=2pt] (-1.63,-1.98) rectangle (1.63,1.43);
  \draw[gray!55,thin] (-1.28,0)--(1.30,0);
  \draw[gray!55,thin] (0,-1.24)--(0,1.14);
  \draw[blue!70!black,line width=1.2pt,smooth] plot[domain=-1.09:0,samples=26] (\x,{0*\x*\x*\x});
  \draw[blue!70!black,line width=1.2pt,smooth] plot[domain=0:1.09,samples=26] (\x,{0*\x*\x*\x});
  \draw[orange!85!black,thick] (-.34,0)--(.34,0);
  \fill[orange!85!black] (0,0) circle (1.6pt);
  \node[font=\scriptsize] at (0,-1.35) {\textcircled{\raisebox{-0.1ex}{\scalebox{0.72}{5}}} 平→平};
\end{scope}
\begin{scope}[shift={(8.21,5.03)},x=.80cm,y=.69cm]
  \draw[gray!35,rounded corners=2pt] (-1.63,-1.98) rectangle (1.63,1.43);
  \draw[gray!55,thin] (-1.28,0)--(1.30,0);
  \draw[gray!55,thin] (0,-1.24)--(0,1.14);
  \draw[blue!70!black,line width=1.2pt,smooth] plot[domain=-1.09:0,samples=26] (\x,{0*\x*\x*\x});
  \draw[blue!70!black,line width=1.2pt,smooth] plot[domain=0:1.09,samples=26] (\x,{-1*\x*\x*\x});
  \draw[orange!85!black,thick] (-.34,0)--(.34,0);
  \fill[orange!85!black] (0,0) circle (1.6pt);
  \node[font=\scriptsize] at (0,-1.35) {\textcircled{\raisebox{-0.1ex}{\scalebox{0.72}{6}}} 平→降};
\end{scope}
\begin{scope}[shift={(1.75,2.51)},x=.80cm,y=.69cm]
  \draw[gray!35,rounded corners=2pt] (-1.63,-1.98) rectangle (1.63,1.43);
  \draw[gray!55,thin] (-1.28,0)--(1.30,0);
  \draw[gray!55,thin] (0,-1.24)--(0,1.14);
  \draw[blue!70!black,line width=1.2pt,smooth] plot[domain=-1.09:0,samples=26] (\x,{-1*\x*\x*\x});
  \draw[blue!70!black,line width=1.2pt,smooth] plot[domain=0:1.09,samples=26] (\x,{1*\x*\x*\x});
  \draw[orange!85!black,thick] (-.34,0)--(.34,0);
  \fill[orange!85!black] (0,0) circle (1.6pt);
  \node[font=\scriptsize] at (0,-1.35) {\textcircled{\raisebox{-0.1ex}{\scalebox{0.72}{7}}} 降→升};
\end{scope}
\begin{scope}[shift={(4.98,2.51)},x=.80cm,y=.69cm]
  \draw[gray!35,rounded corners=2pt] (-1.63,-1.98) rectangle (1.63,1.43);
  \draw[gray!55,thin] (-1.28,0)--(1.30,0);
  \draw[gray!55,thin] (0,-1.24)--(0,1.14);
  \draw[blue!70!black,line width=1.2pt,smooth] plot[domain=-1.09:0,samples=26] (\x,{-1*\x*\x*\x});
  \draw[blue!70!black,line width=1.2pt,smooth] plot[domain=0:1.09,samples=26] (\x,{0*\x*\x*\x});
  \draw[orange!85!black,thick] (-.34,0)--(.34,0);
  \fill[orange!85!black] (0,0) circle (1.6pt);
  \node[font=\scriptsize] at (0,-1.35) {\textcircled{\raisebox{-0.1ex}{\scalebox{0.72}{8}}} 降→平};
\end{scope}
\begin{scope}[shift={(8.21,2.51)},x=.80cm,y=.69cm]
  \draw[gray!35,rounded corners=2pt] (-1.63,-1.98) rectangle (1.63,1.43);
  \draw[gray!55,thin] (-1.28,0)--(1.30,0);
  \draw[gray!55,thin] (0,-1.24)--(0,1.14);
  \draw[blue!70!black,line width=1.2pt,smooth] plot[domain=-1.09:0,samples=26] (\x,{-1*\x*\x*\x});
  \draw[blue!70!black,line width=1.2pt,smooth] plot[domain=0:1.09,samples=26] (\x,{-1*\x*\x*\x});
  \draw[orange!85!black,thick] (-.34,0)--(.34,0);
  \fill[orange!85!black] (0,0) circle (1.6pt);
  \node[font=\scriptsize] at (0,-1.35) {\textcircled{\raisebox{-0.1ex}{\scalebox{0.72}{9}}} 降→降};
\end{scope}
\node[font=\scriptsize,anchor=east] at (.27,7.55) {左升};
\node[font=\scriptsize,anchor=east] at (.27,5.03) {左平};
\node[font=\scriptsize,anchor=east] at (.27,2.51) {左降};
\node[font=\scriptsize] at (5.0,.45) {蓝线：曲线走向；橙线：在中间的水平切线};
\end{tikzpicture}
\end{center}

\textbf{怎样读这九种？}
\begin{itemize}
 \item \textbf{\textcircled{\raisebox{-0.1ex}{\scalebox{0.72}{3}}}、\textcircled{\raisebox{-0.1ex}{\scalebox{0.72}{7}}}是真正的掉头：}先升后降，有局部最高处；先降后升，有局部最低处。
 \item \textbf{\textcircled{\raisebox{-0.1ex}{\scalebox{0.72}{1}}}、\textcircled{\raisebox{-0.1ex}{\scalebox{0.72}{9}}}没有掉头：}左右仍朝同一个方向走，只是在 $c$ 这一点切线水平。
 \item \textbf{\textcircled{\raisebox{-0.1ex}{\scalebox{0.72}{2}}}、\textcircled{\raisebox{-0.1ex}{\scalebox{0.72}{4}}}、\textcircled{\raisebox{-0.1ex}{\scalebox{0.72}{6}}}、\textcircled{\raisebox{-0.1ex}{\scalebox{0.72}{8}}}是进出平路：}可以从上升或下降变平，也可以从平路重新上升或下降。
 \item \textbf{\textcircled{\raisebox{-0.1ex}{\scalebox{0.72}{5}}}两侧都平：}附近整段保持同一高度，当然也满足 $f'(c)=0$。
\end{itemize}

\textbf{一句话：导数为零，看的是这一点；曲线会不会掉头，看的是左右两边。}
\textbf{与罗尔定理的区别：}罗尔定理在两端同高等条件下\textbf{保证至少有一处水平切线}，
却\textbf{没有说找到的每一处水平切线都要掉头}。
\end{explanationbox}
